\documentclass{article}
\usepackage{amsmath}
\usepackage{tikz}
\usetikzlibrary{arrows.meta}

\begin{document}

\begin{figure}[h]
    \centering
    \begin{tikzpicture}[scale=1.5]
        % Draw the left part of the figure
        \draw[thick] (0,0) -- (1,0);
        \filldraw (0,0) circle (2pt);
        \filldraw (1,0) circle (2pt);
        
        % Draw the right part of the figure
        \draw[thick] (2,0) -- (3,0);
        \filldraw (2,0) circle (2pt);
        \filldraw (3,0) circle (2pt);
        
        % Draw the arrows
        \draw[-{Latex[length=3mm]}] (1.5,-0.2) -- (2.5,-0.2);
        
        % Draw the nodes and edges in the right part of the figure
        \node at (2.5,0.5) [circle,fill,inner sep=1.5pt,label=above:1] {};
        \node at (2.5,0) [circle,fill,inner sep=1.5pt,label=below:2] {};
        \node at (2.5,-0.5) [circle,fill,inner sep=1.5pt,label=below:3] {};
        \node at (3.5,0) [circle,fill,inner sep=1.5pt,label=below:5] {};
        \node at (4.5,0) [circle,fill,inner sep=1.5pt,label=below:8] {};
        \node at (5.5,0) [circle,fill,inner sep=1.5pt,label=below:12] {};
        \node at (4.5,0.5) [circle,fill,inner sep=1.5pt,label=above:6] {};
        \node at (4.5,-0.5) [circle,fill,inner sep=1.5pt,label=below:7] {};
        \node at (5.5,0.5) [circle,fill,inner sep=1.5pt,label=above:9] {};
        \node at (5.5,-0.5) [circle,fill,inner sep=1.5pt,label=below:11] {};
        \node at (6.5,0.5) [circle,fill,inner sep=1.5pt,label=above:10] {};
        
        % Draw the edges
        \draw[thick] (2.5,0.5) -- (3.5,0);
        \draw[thick] (2.5,0.5) -- (4.5,0.5);
        \draw[thick] (2.5,0.5) -- (5.5,0.5);
        \draw[thick] (2.5,0.5) -- (6.5,0.5);
        \draw[thick] (2.5,0) -- (3.5,0);
        \draw[thick] (2.5,0) -- (4.5,0);
        \draw[thick] (2.5,0) -- (5.5,0);
        \draw[thick] (2.5,0) -- (6.5,0);
        \draw[thick] (2.5,-0.5) -- (3.5,0);
        \draw[thick] (2.5,-0.5) -- (4.5,-0.5);
        \draw[thick] (2.5,-0.5) -- (5.5,-0.5);
        \draw[thick] (2.5,-0.5) -- (6.5,-0.5);
        
        % Draw the dashed line
        \draw[dashed] (3.5,0) -- (4.5,0);
        
        % Add labels for the caption
        \node at (0.5,0.2) {The solid lines in the figure show one of the non-symmetric spaces, with \( N = 3, L_* = 4 \) constructed in the fashion of Example \ref{ex:nonsym}, together with the ordering of vertices described.};
        \node at (0.5,0) {The dashed line indicates an edge, which can be added to produce another conductively uniform oriented replacement rule with the \( L_* \)-uniform scaling property, where the dashed edge is removable in the sense of Definition \ref{def:removable}.};
        \node at (0.5,-0.2) {In general, any edge which connects vertically aligned points is removable.};
    \end{tikzpicture}
    \caption{}
\end{figure}

\end{document}